\uuid{gYtR}
\exo7id{1940}
\titre{exo7 1940}
\auteur{gineste}
\organisation{exo7}
\datecreate{2001-11-01}
\isIndication{false}
\isCorrection{false}
\chapitre{Série numérique}
\sousChapitre{Convergence absolue}
\cpge{mpsi-mp2i.procedes-sommatoires-discrets.a, pcsi.series-numeriques.a, ptsi.series-numeriques.a, mpsi-mp2i.procedes-sommatoires-discrets.b, pcsi.series-numeriques.b, ptsi.series-numeriques.b, mpsi-mp2i.procedes-sommatoires-discrets.c, pcsi.series-numeriques.c, ptsi.series-numeriques.c, bcpst1.suites-reelles, bcpst1.limites-continuite.b, mp-mpi.series-numeriques-vectorielles.a, pt.series-numeriques.b, mp-mpi.series-numeriques-vectorielles.b, pc.series-numeriques, psi.series-numeriques.a, pt.series-numeriques.a, bcpst2.series-reelles}
\module{Analyse}
\niveau{L1}
\difficulte{}

\contenu{
\texte{
Etudier les séries de termes généraux
}
\begin{enumerate}
    \item \question{$$ u_n=\frac{(-1)^n}{(\ln n) (n^{1/n})} $$}
    \item \question{$$ v_n=\frac{(-1)^n}{\sqrt{n^\alpha + (-1)^n}} \mbox{ o\`u } \alpha > 0 $$}
    \item \question{$$ w_n=\ln(1 + \frac{(-1)^n}{n^\alpha}) \mbox{ o\`u } \alpha > 0 $$}
\end{enumerate}
}
